\documentclass[12pt]{article}
\usepackage[utf8]{inputenc}
\usepackage[T1]{fontenc}
\usepackage{hyperref}
\title{Savi Language Resources for K-12 Curriculum Development}
\author{Compiled by Claude Sonnet 5.5 (claude-sonnet-5-5) for Andr\'as Kornai}
\date{}
\begin{document}
\maketitle

\section{Introduction}
Savi (also Sawi, Sauji; ISO 639-3 sdg; Glottocode savi1242) is an endangered Indo-Aryan (Dardic, Shina-cluster) language of northeastern Afghanistan, centred on the village of Sau on the Kunar River south of Arandu, near the Pakistani border. Published speaker estimates vary widely (roughly 4,000 to 12,000), and refugee dispersal since the Afghan wars makes current numbers uncertain. Little published material exists. Searches found no school materials, primers or dictionaries.

\section{Primary Educational Resources}
No primers, textbooks or school programmes for Savi were found in search results, and the language is not known to be taught in schools.

\section{Linguistic and Reference Materials}
\begin{itemize}
\item Wikipedia overview of Sawi (Dardic), with classification, location and speaker estimates: \url{https://en.wikipedia.org/wiki/Sawi_language_(Dardic)}
\item Endangered Languages Project entry for Savi: \url{https://endangeredlanguages.com/lang/5362}
\item Numerals 1--2000 in Savi (attributed on the page to Henrik Liljegren): \url{https://lingweb.eva.mpg.de/channumerals/Savi.htm}
\item Wiktionary has a few Savi entries. Per the search summary, they cite a grammar sketch of Sauji by Nina Knobloch (2020, Stockholm University). A direct search for that sketch found nothing, so its existence is unverified here. \url{https://en.wiktionary.com/wiki/Sawi}
\end{itemize}

\section{Religious and Community Resources}
No scripture translations or community organisations for Savi were found.

\section{Digital and Audio Resources}
No audio, video or software resources specific to Savi were found. The numerals table above is the only lexical dataset located.

\section{Curriculum Development Context}
Savi has no known orthography, literacy materials or school use. Any curriculum would first need community consultation, documentation and an orthography. The sources above are suitable starting points for linguistic documentation only. Dari and Pashto are the likely languages of schooling in the area.
\end{document}